\documentclass[conference]{IEEEtran}
\IEEEoverridecommandlockouts

\usepackage{cite}
\usepackage{amsmath,amssymb,amsfonts}
\usepackage{graphicx}
\usepackage{booktabs}
\usepackage{multirow}
\usepackage{xcolor}
\usepackage{url}

% TODO notes while drafting -- remove before submission
\newcommand{\todo}[1]{\textcolor{red}{[TODO: #1]}}

\begin{document}

\title{Classification of Marine Accident Investigation Reports\\
Using Transformer-Based Language Models}

\author{
\IEEEauthorblockN{Mohammad Mahdi Arab}
\IEEEauthorblockA{\textit{Faculty of Computer Engineering} \\
\textit{Shahrood University of Technology}\\
Shahrood, Iran \\
email@example.com}
\and
\IEEEauthorblockN{Supervisor Name}
\IEEEauthorblockA{\textit{Faculty of Computer Engineering} \\
\textit{Shahrood University of Technology}\\
Shahrood, Iran \\
email@example.com}
}

\maketitle

\begin{abstract}
\todo{150--250 words. Problem -> gap -> method -> dataset -> key quantitative result -> implication. Write this LAST.}
\end{abstract}

\begin{IEEEkeywords}
maritime safety, accident reports, text classification, BERT, natural language processing
\end{IEEEkeywords}

% ------------------------------------------------------------
\section{Introduction}
\todo{Why marine incident reports matter; why manual classification does not scale; the gap; your contribution.}

Contributions of this paper:
\begin{itemize}
  \item \todo{Contribution 1 -- e.g., a labeled dataset built from MAIB reports}
  \item \todo{Contribution 2 -- e.g., systematic comparison of models}
  \item \todo{Contribution 3 -- e.g., error analysis / practical insight}
\end{itemize}

% ------------------------------------------------------------
\section{Related Work}
\subsection{Text Mining in Maritime Safety}
\todo{Prior work on accident report analysis.}

\subsection{Transformer Models for Domain-Specific Classification}
BERT~\cite{devlin2019bert} \todo{BERT and domain adaptation.}

% ------------------------------------------------------------
\section{Dataset}
\todo{Source (MAIB), collection, preprocessing, label definitions, class distribution.}

\begin{table}[htbp]
\caption{Dataset Statistics}
\label{tab:dataset}
\centering
\begin{tabular}{lrr}
\toprule
Class & \# Reports & \% \\
\midrule
Class A & 0 & 0.0 \\
Class B & 0 & 0.0 \\
\midrule
Total & 0 & 100.0 \\
\bottomrule
\end{tabular}
\end{table}

% ------------------------------------------------------------
\section{Methodology}
\todo{Model architecture, input handling (long reports vs.\ 512-token limit), fine-tuning procedure.}

% ------------------------------------------------------------
\section{Experimental Setup}
\todo{Train/val/test split, baselines, hyperparameters, hardware, evaluation metrics (macro-F1 for imbalance).}

% ------------------------------------------------------------
\section{Results and Discussion}

\begin{table}[htbp]
\caption{Classification Performance}
\label{tab:results}
\centering
\begin{tabular}{lccc}
\toprule
Model & Precision & Recall & Macro-F1 \\
\midrule
Baseline (TF-IDF + SVM) & -- & -- & -- \\
BERT & -- & -- & -- \\
\bottomrule
\end{tabular}
\end{table}

As shown in Table~\ref{tab:results}, \todo{interpretation}.

% Example figure:
% \begin{figure}[htbp]
% \centerline{\includegraphics[width=\columnwidth]{figures/confusion_matrix.pdf}}
% \caption{Confusion matrix of the best model.}
% \label{fig:cm}
% \end{figure}

% ------------------------------------------------------------
\section{Conclusion}
\todo{Summary, limitations, future work.}

\section*{Acknowledgment}
\todo{Optional.}

\bibliographystyle{IEEEtran}
\bibliography{refs}

\end{document}
